\section{Evaluation Design}\label{sec:evaluation_plan}

This article \emph{proposes} a reference architecture and \emph{instantiates} it as a proof of concept; it does not claim a completed empirical validation. Accordingly, this section presents an evaluation \emph{plan} rather than results: for each research question stated in \cref{sec:introduction} it specifies a study design, the corpus on which the study will run, the baselines against which the pipeline will be compared, the metrics that will be collected, and the analysis that will be applied to them. The intent is to make the architecture falsifiable---to state in advance what observations would count as evidence for, and against, each of the claims the pipeline makes---and to fix the operational definitions of the metrics so that later work, by these authors or others, can execute the plan without reinterpreting it. Where the proof-of-concept instantiation of \cref{sec:instantiation} already exercises part of a study, this is noted as evidence of feasibility; full effect-size estimation on external production codebases, together with an industrial case study, is deferred to future work for the reasons set out in \cref{subsec:eval_industrial}.

Two design principles govern the whole plan. First, every metric must be \emph{tool-derived and objective} wherever possible: the pipeline records provenance, assurance level, counterexamples, iteration counts, and verification outcomes as a by-product of its normal operation, so the primary measurements are extracted from machine logs rather than from human judgement. Where human judgement is unavoidable---confirming that a candidate bug is real, or that a ratified example reflects true intent---the protocol fixes the adjudication procedure in advance and reports inter-rater agreement. Second, the trusted verifier is never permitted to grade its own inputs: the OpenJML extended static checker (ESC) discharged by the Z3 SMT solver, together with the custom fork adding \texttt{define-fun-rec} for \verb|\sum|, \verb|\product|, and \verb|\num_of|, is a fixed instrument reused unchanged from the authors' prior work, and all measurements that depend on it are reported relative to that fixed oracle.

\subsection{Corpus}\label{subsec:eval_corpus}

The primary corpus is Apache Commons Lang, chosen because it is the rare real-world artefact that exposes all three intent streams the brownfield workflow depends upon. Its source supplies the \emph{code} stream; the public Apache JIRA project (issues of the form \texttt{LANG-\#\#\#\#}) supplies the \emph{ticket} stream, including reported defects and discussed edge cases; and its rich Javadoc supplies the \emph{documentation} stream. The library is already integrated into the authors' experimental harness from prior work~\cite{edmonds_article3,edmonds_article4}, so the inference instrument, the converter, and the Dockerized verification pipeline run on it without modification, and the corpus is large enough to exercise propagation and characterisation at scale while remaining small enough to permit per-clause human adjudication on a curated subset. Studies that require per-method ground truth are scoped to the \texttt{StringUtils} and \texttt{NumberUtils} classes, where code, tickets, and Javadoc are densest and where the method classes the inferrer already handles well---null-safety, array bounds, simple arithmetic postconditions, and the loop patterns (\textsc{max}/\textsc{min}, linear search, sum) for which it emits invariants---are best represented. Two supplementary corpora support external-validity checks: a held-out set of public GitHub Java projects not seen during heuristic development, used to test whether characterisation and propagation generalise beyond the tuning corpus, and the SWE-bench issue set~\cite{jimenez2024swebench} as a point of comparison for the agentic synthesis baselines defined below. The held-out corpora are used only for feasibility and generalisation observations in this article; controlled measurement on them is future work.

\subsection{RQ1: Ratification effort and the ratification gap}\label{subsec:eval_rq1}

The first study asks whether the \emph{ratify-behaviour-not-logic} mechanism (\cref{sec:intent_layer}) lets a human own the contract at acceptable cost, given that the human chose NL precisely because they cannot or will not read formal text. The design is a within-subjects task study in which participants elaborate a fixed set of method specifications under two conditions: a \emph{formal-text} baseline, in which the participant confirms the JML contract directly, and the \emph{example-pinned} condition, in which the agent renders the contract as concrete checkable examples (``given \verb|x=-1| this contract says \verb|result=0|; correct?'') and the participant ratifies behaviour rather than logic. The corpus is the curated \texttt{StringUtils}/\texttt{NumberUtils} subset, ordered to balance difficulty across conditions.

The primary metric is \emph{ratification effort}, defined as obligations confirmed per minute of participant time, measured from interaction logs; the secondary metric is the \emph{fraction auto-defaulted}, the proportion of obligations resolved by inherited default contracts (\cref{sec:intent_layer}) without participant adjudication, since the architecture's claim is that defaults absorb the bulk of routine intent and leave only departures for the human. A faithfulness metric guards against the failure mode in which examples are easy but wrong: each ratified example becomes a \emph{pinned oracle}, and the study records how many pinned oracles are later contradicted by a participant on re-presentation, a direct measure of whether behavioural ratification produces stable intent. Analysis pairs the two conditions per participant and reports the effect on effort with a non-parametric paired test and confidence intervals, alongside the auto-default fraction as a descriptive distribution. The hypothesis is that the example-pinned condition raises obligations-per-minute and that inherited defaults account for a large share of obligations, narrowing the ratification gap without inflating the contradiction rate.

\subsection{RQ2: Decision surfacing and adversarial soundness}\label{subsec:eval_rq2}

The second study evaluates the two elaboration-layer guards: Socratic decision surfacing and adversarial elaboration (\cref{sec:intent_layer}). It uses a seeded-decision design. For a set of methods drawn from the corpus, the authors construct ground-truth catalogues of boundary decisions---null and empty handling, integer overflow, ordering, duplicates, concurrency, and the error-versus-exception choice---that a correct elaboration must resolve. The \emph{gap/boundary detection rate} is then the fraction of seeded decisions the agent surfaces to the human \emph{before} code is synthesized, versus those that escape surfacing and would otherwise be discovered only at or after verification; escaped decisions are counted by replaying synthesis and observing where an unsurfaced choice manifests as a counterexample or a divergence from the pinned oracles.

Adversarial elaboration is evaluated on its ability to catch \emph{spec weakening}. The study seeds contracts with known weakenings---most importantly the canonical case in which ``sort'' is degraded to ``is a permutation'', which the identity function trivially satisfies---and measures the \emph{spec-weakening catch rate} of the second, independent agent that searches for scenarios consistent with the NL intent but violating the proposed contract. Two further detectors are measured against seeded faults: the \emph{vacuity rate}, the fraction of contracts flagged because a postcondition is trivially true or a precondition is unsatisfiable, repurposing the authors' prior ``\texttt{ensures true}'' punt detection as a weakening signal; and the resulting precision and recall over the seeded weakening and vacuity faults. Baselines are a single-agent elaboration with no adversarial pass and a syntactic vacuity check alone. The analysis reports detection rate, catch rate, and precision/recall with binomial confidence intervals, and the hypothesis is that the adversarial pass and vacuity detector together raise weakening recall substantially over the single-agent baseline at acceptable precision.

\subsection{RQ3: Provenance, assurance, and independent authority}\label{subsec:eval_rq3}

The third study asks whether the contract that code is verified against is honestly and auditably constituted---that is, whether the independent-authority constraint and the provenance/assurance tagging make the ``verified'' claim defensible (\cref{sec:architecture}). The unit of observation is the obligation at merge time. For every merged contract the pipeline emits each obligation's provenance tag (human-ratified, inherited-default, inferred-unconfirmed, example-pinned, ticket-stated, doc-stated, or code-inferred) and assurance tag (proved, bounded-checked, tested, or unchecked). The primary reported artefact is the \emph{provenance distribution at merge time}, a contingency table over provenance $\times$ assurance, which makes visible what fraction of a merged contract rests on human-ratified, proved obligations versus inferred-unconfirmed, merely tested ones. A healthy distribution concentrates the load-bearing, oracle-defining obligations in the human-ratified and example-pinned rows; a distribution dominated by code-inferred, unconfirmed clauses signals a contract that is characterisation, not specification (\cref{subsec:eval_rq5}).

Two integrity checks accompany the distribution. An \emph{authorship-independence audit} verifies, from the orchestration log, that no obligation in the merged contract carries an author identity matching the implementer role that produced the code it verifies; any violation is a circularity defect and is reported as a count. A \emph{degradation-accounting} check verifies that every obligation that could not be proved is recorded at a strictly weaker assurance level (bounded-checked or tested) rather than silently dropped, by reconciling the set of attempted obligations against the set of merged obligations. The analysis is descriptive: the contingency table, the independence-violation count (target zero), and the no-silent-drop reconciliation. \Cref{tab:eval_summary} summarises the metric for each research question.

\subsection{RQ4: The synthesis--verification (CEGIS) loop}\label{subsec:eval_rq4}

The fourth study evaluates the counterexample-guided synthesis loop (\cref{sec:verification_loop}). For each target method, the agent emits code, OpenJML verifies the code against the fixed contract, and on failure the concrete counterexample witness is returned to the synthesis agent, which regenerates, capped at $N$ iterations; on persistent failure the obligation is degraded per the tiered-assurance policy and reported. The corpus is stratified by \emph{method class}: null-safety, array bounds, simple arithmetic postconditions, and the inferrer-supported loop patterns, since the architecture's claim is class-dependent convergence rather than uniform success.

The primary metrics are the \emph{success rate by method class}---the fraction of methods for which a contract-satisfying implementation is found within the iteration cap---and the \emph{distribution of loop iterations to convergence}, reported per class so that easy classes (null-safety, bounds) are not averaged together with hard ones (multiplicative arithmetic, rich preconditions). To isolate the contribution of the counterexample channel, the principal baseline replaces the concrete witness with a test-failure signal (``an input failed'') of the kind the agentic baselines of \cref{subsec:rw0} rely on~\cite{shinn2023reflexion,chen2024selfdebug}; a second baseline is open-loop sampling with no verification gate, evaluated against a strong test suite~\cite{liu2023evalplus} to expose the plausible-but-incorrect outputs the closed loop is designed to reject. Milestone M1 of the instantiation---the synthesis loop run against a fixed hand-written contract on an already-green method class---supplies an initial feasibility data point for this study and de-risks the loop before contract generation is introduced. The analysis compares success rate and median iterations across conditions per class, with the hypothesis that the counterexample channel converges in fewer iterations and at higher success rate than the test-failure channel, and that the gap widens on harder classes where a concrete witness carries more diagnostic information than ``some test failed''.

\subsection{RQ5: Brownfield characterisation and candidate-bug yield}\label{subsec:eval_rq5}

The fifth study evaluates the startup/characterisation phase and the multi-source intent triangulation it performs (\cref{sec:brownfield}). The phase ingests code, JIRA tickets, and Javadoc; infers a draft contract from the code; extracts intent claims from tickets and docs; and reconciles them into a candidate contract with per-clause provenance, flagging disagreements between streams. The study runs this phase over the \texttt{StringUtils}/\texttt{NumberUtils} subset and treats the disagreements as the output of interest, in keeping with the principle that this phase cannot be gated on zero errors: finding code-versus-intent violations is the success outcome, not a failure.

The primary metric is \emph{characterisation discrepancy yield}: the number of candidate bugs surfaced as disagreements---either code violating its own inferred pattern (an over-strong inferred clause) or code disagreeing with ticket- or doc-derived intent---partitioned into \emph{confirmed} and \emph{false} by a fixed adjudication protocol in which two raters inspect each candidate against the originating ticket or Javadoc and the source, with disagreements resolved by discussion and inter-rater agreement reported. Because the corpus exposes historical \texttt{LANG-\#\#\#\#} defects, a recall proxy is available: the fraction of historically documented edge-case defects in the subset that the discrepancy stream re-surfaces. The verification-as-characterisation caveat is measured explicitly rather than asserted: the study reports the fraction of merged clauses that are code-derived (and therefore nearly circular, proving faithful description and toolchain sanity but not correctness) separately from the fraction grounded in ticket/doc intent plus human ratification, making precise what the verification does and does not establish. A secondary measurement records \emph{traceability noise}---the precision of the heuristic and retrieval-augmented links between tickets/docs and code~\cite{antoniol2002recovering,borg2014recovering}---since the credibility of the triangulation rests on those links and they are known to be noisy. The expected outcome is a non-trivial confirmed candidate-bug yield on a mature, well-tested library, demonstrating that disagreement detection finds real defects rather than merely re-describing the code.

\subsection{RQ6: Compositional change propagation and version compatibility}\label{subsec:eval_rq6}

The sixth study evaluates the lifecycle capability that links this article to the thesis core: because contracts compose, a callee-spec change recomputes caller proof obligations and flags which callers break without running them, so that an internal refactor and a dependency or library version bump are the same operation (\cref{sec:propagation}). This promotes the authors' prior compositional refinement~\cite{edmonds_article3} from a one-off inference result to a continuous blast-radius analysis. The design uses real version pairs from Commons Lang and from the broader Maven ecosystem, where behavioural breaking changes between releases have been catalogued in prior empirical work~\cite{ochoa2022breaking,mostafa2017behavioral,jayasuriya2024understanding}. For a callee whose contract changes between versions, the pipeline computes the set of call sites whose recomputed obligations no longer discharge---the predicted breaking callers---and this set is compared against ground truth derived from the catalogued behavioural incompatibilities and from differential execution of the affected clients.

The primary metrics are \emph{propagation precision and recall} over flagged breaking callers: precision penalizes false alarms that would erode developer trust, and recall penalizes missed breakages that defeat the purpose of the analysis. Baselines are signature-level breaking-change detectors~\cite{brito2018apidiff} and semantic-versioning declarations~\cite{preston2013semver,raemaekers2017semver}, neither of which reasons about behavioural contracts, so the hypothesis is that contract-level propagation recovers behavioural breakages these miss while not flagging benign signature-compatible changes. A cost metric records the analysis time per call site, since the value of computing the blast radius ``without running them'' depends on it being cheaper than differential execution. The analysis reports precision, recall, and their harmonic mean against each baseline, broken down by whether the change is an internal refactor or a cross-version library bump, to substantiate the claim that the two are handled by one mechanism.

\subsection{Cross-cutting metric: verification cost and graceful degradation}\label{subsec:eval_cost}

Every study above depends on the trusted oracle, whose cost and decidability limits are real and known---Z3 times out on multiplicative reasoning, array theory, and rich preconditions, and not all intent is formally specifiable. The plan therefore collects, across all studies, the \emph{verification cost} (wall-clock discharge time per obligation, under the settled strict configuration of code-math \texttt{safe} and spec-math \texttt{bigint}) and the \emph{timeout rate} (the fraction of obligations on which the solver does not terminate within the budget), reported per method class so that the cost profile of the architecture is transparent. These two measurements feed the tiered-assurance accounting of \cref{subsec:eval_rq3}: a timeout is not a verification failure but a trigger to degrade full proof to bounded check to test, recording the assurance level. Reporting cost and timeout rate alongside the per-RQ metrics ensures that any feasibility claim is qualified by the honest statement that verification does not always terminate, and that the architecture's graceful-degradation response is itself measured rather than assumed.

\begin{table}[t]
\centering
\caption{Evaluation plan summary: primary metric and baseline per research question.}
\label{tab:eval_summary}
\small
\begin{tabular}{@{}llll@{}}
\toprule
RQ & Mechanism evaluated & Primary metric & Principal baseline \\
\midrule
RQ1 & Behavioural ratification        & Obligations confirmed/min; auto-default fraction & Formal-text confirmation \\
RQ2 & Surfacing + adversarial pass    & Boundary detection rate; weakening catch rate    & Single-agent elaboration \\
RQ3 & Provenance + independent authority & Provenance$\times$assurance distribution        & --- (auditability check) \\
RQ4 & CEGIS synthesis loop            & Success rate \& iterations by class              & Test-failure feedback signal \\
RQ5 & Brownfield characterisation     & Confirmed candidate-bug yield                    & Code-only inference \\
RQ6 & Compositional propagation       & Breaking-caller precision/recall                 & Signature diff; semver \\
\bottomrule
\end{tabular}
\end{table}

\subsection{Planned industrial case study (future work)}\label{subsec:eval_industrial}

The studies above can be executed on public corpora by the authors alone, but the architecture's central claims concern a \emph{lifecycle}---how teams elaborate, ratify, synthesize, and maintain contracts over time---which an offline corpus cannot fully exercise. The plan therefore includes an industrial case study, framed explicitly as future work, conducted as \emph{action research} within a partner organisation that adopts the pipeline on a live codebase. The methodological commitment is that the case study measures \emph{objective, tool-derived metrics}---the same ratification-effort, provenance-distribution, candidate-bug-yield, CEGIS-convergence, and propagation-precision quantities defined above, extracted from the running pipeline's logs---rather than self-report instruments such as developer-satisfaction surveys administered by a participant's manager, whose validity is compromised by the power relationship. Human subjects are involved only to the extent of consenting to log analysis and participating in the ratification dialogue, and the study will proceed under approval from a Human Research Ethics Committee (HREC) with informed consent and data-handling protocols fixed in advance. The retroactive-then-prospective loop reordering of \cref{sec:brownfield} structures the case study: the first phase resolves recovered discrepancies on existing code to establish a trusted baseline and to let the tool earn trust before it changes anything, and only subsequent phases apply the pipeline to new features. Consistent with the limitations enumerated throughout this article, the case study aims to establish \emph{feasibility and realism in situ}---that the pipeline functions, that its metrics are obtainable, and that practitioners can operate it---rather than a controlled effect size, which would require a larger multi-site programme beyond the scope of this work.
